\documentclass[11pt,a4paper]{article}
\usepackage[utf8]{inputenc}
\usepackage[T1]{fontenc}
\usepackage{geometry}
\usepackage{xcolor}
\usepackage{hyperref}
\usepackage{longtable}
\usepackage{booktabs}
\usepackage{enumitem}
\usepackage{fancyhdr}
\usepackage{listings}
\usepackage{titlesec}
\usepackage{parskip}
\usepackage{amsmath,amssymb,amsthm}
\usepackage{multirow}

\geometry{margin=1in}

\definecolor{luxblue}{RGB}{30,64,175}
\definecolor{codebg}{RGB}{245,245,245}
\definecolor{darkgray}{RGB}{60,60,60}
\definecolor{statusdone}{RGB}{0,128,64}
\definecolor{statusdev}{RGB}{200,120,0}
\definecolor{statustodo}{RGB}{180,40,40}

\hypersetup{
  colorlinks=true,
  linkcolor=darkgray,
  urlcolor=luxblue,
  citecolor=darkgray,
}

\lstset{
  basicstyle=\small\ttfamily,
  backgroundcolor=\color{codebg},
  frame=single,
  framerule=0pt,
  breaklines=true,
  columns=fullflexible,
  literate={\_}{\_}1,
}

\titleformat{\section}{\Large\bfseries}{}{0em}{}[\vspace{-0.5em}\rule{\textwidth}{0.4pt}]
\titleformat{\subsection}{\large\bfseries}{}{0em}{}
\titleformat{\subsubsection}{\normalsize\bfseries}{}{0em}{}

\pagestyle{fancy}
\fancyhf{}
\fancyhead[L]{\small\textit{Lux --- GPU-Native Encrypted Mempool}}
\fancyhead[R]{\small\textit{Lux Industries, Inc.}}
\fancyfoot[C]{\small\thepage}

\newtheorem{theorem}{Theorem}
\newtheorem{lemma}{Lemma}
\newtheorem{definition}{Definition}
\newtheorem{proposition}{Proposition}

\newcommand{\done}{\textcolor{statusdone}{\textbf{SHIPPED}}}
\newcommand{\indev}{\textcolor{statusdev}{\textbf{IN BUILD}}}
\newcommand{\todesign}{\textcolor{statustodo}{\textbf{SPEC ONLY}}}

\begin{document}

\begin{center}
{\Huge\bfseries A GPU-Native Encrypted Mempool}\\[0.5em]
{\Large Threshold TFHE with Deferred Decryption at Block Construction}\\[1.5em]
{\large Lux Industries, Inc.}\\
\texttt{lux.network}\\[0.5em]
{\small \today}
\end{center}

\vspace{1em}

\begin{abstract}
\noindent
The Lux protocol stack ships a fully homomorphic encryption (FHE)
precompile at \texttt{0x0700\dots} on every EVM-compatible chain in
the family. This paper specifies an end-to-end encrypted-mempool
protocol layered on the existing precompile: clients encrypt order
intents under a network-wide threshold TFHE public key; validators
admit ciphertexts to the mempool on signature, fee, and a NIZK
well-formedness proof; the block proposer threshold-decrypts the
batch at block construction time and feeds plaintext intents to the
matching engine. The entire pipeline runs natively on GPU --- client
encryption via WebGPU, CUDA, or Metal Performance Shaders depending
on the consumer; validator partial-decryption on the same H100/H200
fleet already provisioned for the TEE attestation precompile
(\texttt{0x7201}). The design eliminates classical maximal
extractable value (MEV) on the protected surface --- sandwich and
priority-front-running attacks become computationally infeasible
under TFHE semantic security --- without requiring users to opt in
per trade, choose a builder, or route through a private relay. We
give the protocol, the GPU-native pipeline, the security argument,
and a punch list of remaining work against the existing upstream
modules \texttt{luxfi/fhe}, \texttt{luxfi/threshold},
\texttt{luxfi/node}, \texttt{luxfi/consensus}, \texttt{luxfi/dex},
\texttt{luxfi/vm}, and \texttt{luxfi/bridge}.
\end{abstract}

\tableofcontents
\vspace{1em}

% ================================================================
\section{Introduction}
% ================================================================

A pending transaction on a public mempool is a public broadcast:
anyone can read it, compute the effect of executing it, and
synthesise transactions that profit at the submitter's expense ---
sandwich attacks, front-running of liquidations, time-bandit
re-orgs, priority-gas auctions. The Lux stack's existing answer to
this is multifaceted: validator sets that may be configured
permissionless or permissioned per chain, an FHE precompile
already activated at \texttt{0x0700\dots} for encrypted state, and
a threshold-signature substrate (\texttt{luxfi/threshold}) that
gates cross-chain and consensus signatures behind $M$-of-$N$ share
holders.

What is missing from the surface today is a single ingredient: the
glue that combines those pieces into an end-to-end encrypted
mempool. The FHE precompile can decrypt; the mempool cannot
currently admit ciphertexts; the proposer cannot currently invoke a
batch decryption at block construction. This paper specifies that
glue, GPU-native, and tracks the build against existing upstream
tickets.

% ================================================================
\section{Threat Model}
% ================================================================

\subsection{Attacker Capabilities}

We assume an attacker with the following standard MEV capabilities:

\begin{enumerate}[nosep,label=\textbf{C\arabic*.}]
\item Observe every pending transaction in any public mempool.
\item Submit transactions of arbitrary size at arbitrary fee.
\item Collude with a subset of validators bounded below the
  byzantine threshold of the consensus protocol.
\item Operate flash-loan markets to access arbitrary short-duration
  capital.
\item Observe and re-org public chains within the network's
  finality depth.
\end{enumerate}

\noindent We do \emph{not} assume the attacker can break standard
cryptographic primitives (ECDSA, BLS, ML-KEM-768, ML-DSA-65, TFHE
under Ring-LWE hardness) or break the TEE attestation primitives
that the platform's compliance sidecars rely on.

% ================================================================
\section{Protocol}
% ================================================================

\subsection{Pipeline}

\begin{enumerate}[nosep,label=\textbf{P\arabic*.}]
\item Client encrypts the order tuple
  $(\text{symbol}, \text{side}, \text{quantity}, \text{limit})$
  under the network's threshold TFHE public key. Attaches a NIZK
  proof of well-formedness.
\item Encrypted intent enters the mempool. Validators see a
  ciphertext, the fee bid, a signature over the
  ciphertext-hash, and the NIZK proof. They do not see content.
\item Validators order pending intents by arrival time (FIFO),
  per-account partitioned to prevent self-flooding.
\item At block construction the proposer batches $N$ encrypted
  intents and invokes the FHE precompile to threshold-decrypt
  the batch. The proposer's GPU sidecar produces its partial
  decryption; the proposer collects $M$-of-$N$ partials from peer
  sidecars over the intra-cluster mesh; the precompile combines
  them.
\item The DEX matching engine executes the plaintext batch
  against orderbook state in the chosen FIFO order, within the
  2~s block window.
\end{enumerate}

\subsection{Wire format}

A new transaction type (\texttt{EncryptedPayloadTx}) carries:

\begin{itemize}[nosep]
\item \texttt{ciphertext}: SnP-packed TFHE ciphertext, $\sim$10--30~KB.
\item \texttt{nizk\_proof}: well-formedness proof; constant-size,
  $\sim$1~KB.
\item Plaintext metadata: signature, fee bid, sender, expiry,
  bootstrap-budget bid.
\end{itemize}

\subsection{Reference precompile address space}

The Lux EVM reserves \texttt{0x0700\dots} for privacy and encryption
precompiles. The FHE precompile occupies the first slot in that
range, with adjacent slots for ACL, input verifier, and gateway
contracts. The encrypted mempool reuses this same precompile; no
new address is required.

% ================================================================
\section{GPU-Native Pipeline}
% ================================================================

The cost driver in this protocol is the TFHE bootstrap operation
performed inside the precompile during step P4. On modern hardware
this is the difference between an unbuildable system (CPU TFHE
bootstrap at $\sim$10~ms/gate) and a buildable one (H100/H200 GPU
TFHE bootstrap at $\sim$10--50~\textmu s/gate, batched).

\subsection{Client-side encryption}

The client SDK exposes a single entry point,
\texttt{encryptIntent(\{symbol, side, qty, limit\})}, returning
$(\text{ciphertext}, \text{NIZK\_proof})$. Implementation chooses
the backend at runtime via capability detection:

\begin{longtable}{p{3.4cm}p{5.0cm}p{6.2cm}}
\toprule
\textbf{Tier} & \textbf{Backend} & \textbf{Audience} \\
\midrule
Native &
  \texttt{luxfi/lattice/v7/gpu} CUDA on NVIDIA, Metal on Apple &
  Market maker bots, institutional traders, validators \\
Browser GPU &
  WebGPU compute shaders against the same kernels &
  Retail web users on modern browsers \\
Edge CPU fallback &
  Native Go fall-through &
  Old browsers, last-resort path \\
\bottomrule
\end{longtable}

The NIZK proof keeps mempool admission cheap: validators verify the
proof against the precompile without ever decrypting. Malformed
ciphertexts cannot reach block construction.

\subsection{Validator pipeline}

Per-validator architecture:

\begin{itemize}[nosep]
\item \textbf{Mempool extension}: admit on signature + fee + NIZK
  + bootstrap-budget bid, no decryption. Per-account FIFO inside
  each partition.
\item \textbf{Block proposer hook} (every 2~s): proposer collects
  $N$ encrypted intents into a batch.
\item \textbf{Threshold partial-decrypt service}: a GPU sidecar
  running the \texttt{luxfi/fhe} threshold partial-decrypt
  surface, exposed via ZAP envelope (LP-200). Holds one share of
  the network FHE secret key. Computes its partial decryption of
  the batch in one bootstrapped pass.
\item \textbf{Cross-validator combine}: proposer collects
  $M$-of-$N$ partial decryptions from peer sidecars over the
  intra-cluster mesh. Combine is fast linear algebra on the
  proposer's GPU.
\item \textbf{Plaintext feed}: handed to the DEX matching engine
  for execution.
\end{itemize}

\subsection{Why threshold TFHE on GPU}

\begin{longtable}{p{3.8cm}p{4.2cm}p{2.4cm}p{4.4cm}}
\toprule
\textbf{Option} & \textbf{Security} & \textbf{GPU fit} & \textbf{Verdict} \\
\midrule
Single-key TFHE &
  One holder of FHE secret key &
  Trivial GPU port &
  Single point of compromise. \textbf{Reject.} \\
TEE-only &
  TEE-attested boundary &
  Existing \texttt{0x7201} precompile &
  Strong but trusts hardware vendor. \\
Threshold TFHE on GPU &
  $M$-of-$N$ secret-sharing of FHE key &
  \texttt{luxfi/lattice/v7/gpu} + thin partial-decrypt &
  No single point of compromise. Throughput matches block
  time. \textbf{Recommended.} \\
Hybrid: Threshold + TEE &
  TEE attests each share's correct use &
  Same GPU; defense-in-depth &
  Marginal incremental win. \textbf{Add later.} \\
\bottomrule
\end{longtable}

\subsection{Block-level bootstrap budget}

Each block carries a TFHE bootstrap meter,
$\texttt{tfhe\_ops\_used} / \texttt{tfhe\_ops\_max}$, analogous to
EVM gas. This prevents one large encrypted order from starving the
block by forcing an unbounded bootstrap chain. The meter is enforced
inside the precompile dispatch and surfaced as a transaction-level
fee bid: a higher bid lets a transaction consume more bootstrap
budget within its slot.

% ================================================================
\section{Security}
% ================================================================

\begin{theorem}[No sandwich on the encrypted mempool]
\label{thm:sandwich}
Under the protocol of \S3 with a correctly-generated threshold
TFHE key, an honest majority of $M$-of-$N$ share holders, and
TFHE-CCA1 security under Ring-LWE hardness, an attacker cannot
execute a sandwich attack against a target intent.
\end{theorem}

\begin{proof}[Proof sketch]
A sandwich requires the attacker to (a)~identify the victim intent
in the mempool, (b)~determine its direction and size, and (c)~place
front-running and back-running intents that bracket it. Steps (a)
and (b) require distinguishing or recovering plaintext from
ciphertexts under the threshold key; under TFHE semantic security
this is computationally infeasible. Step (c) requires the
attacker's intents to land immediately before and after the
victim's; mempool ordering is FIFO and the validator does not
learn which ciphertext is the victim until the block is built. The
attacker has no targeting primitive. \qed
\end{proof}

\begin{theorem}[Negative expected value of blind front-running]
\label{thm:blind-frontrun}
An attacker who submits speculative orders without information
about pending intents has expected P\&L
\[
\mathbb{E}[\textit{P\&L}_{\textit{blind}}] = -f_{\textit{taker}}
\cdot V_{\textit{blind}}
\]
where $f_{\textit{taker}}$ is the taker fee and
$V_{\textit{blind}}$ is the speculative-order notional.
\end{theorem}

\begin{proof}
With no information advantage the directional bet is a zero-mean
random variable. The attacker pays taker fees on every speculative
intent. Expected P\&L is strictly negative. \qed
\end{proof}

\begin{theorem}[Quorum failure preserves liveness]
\label{thm:quorum-fail}
If $M$-of-$N$ partial decryptions are unavailable at block
construction --- e.g., a partition isolates a validator subset or
$N-M+1$ sidecars are simultaneously down --- the proposer falls
back to plaintext FIFO without dropping the block.
\end{theorem}

\begin{proof}[Proof sketch]
The proposer hook of \S3.4 has a configurable budget
($\Delta t_{\max}$, default 500~ms). If the threshold-decrypt
service does not complete in $\Delta t_{\max}$, the proposer
discards the partial decryptions received and falls back to plain
FIFO. Liveness is preserved at the cost of one block of MEV
exposure on the encrypted-mempool path. Subsequent blocks return
to the protected path as soon as the quorum is restored. \qed
\end{proof}

% ================================================================
\section{Upstream Module Map}
\label{sec:upstream-map}
% ================================================================

\begin{longtable}{p{3.0cm}p{6.0cm}p{2.0cm}p{4.0cm}}
\toprule
\textbf{Module} & \textbf{Role} & \textbf{Status} & \textbf{Reference issue} \\
\midrule
\texttt{luxfi/precompile/fhe} &
  FHE precompile at \texttt{0x0700\dots} &
  \done & --- \\
\texttt{luxfi/fhe} &
  TFHE library (native Go, Ring-LWE, GPU NTT) &
  \done (CPU + GPU NTT) &
  Bootstrap-on-GPU follow-on: \texttt{luxfi/fhe\#4} \\
\texttt{luxfi/threshold} &
  Threshold protocols (CGGMP21, FROST, TFHE keygen) &
  \indev (TFHE distributed decrypt stubbed) &
  Critical path: \texttt{luxfi/threshold\#20} \\
\texttt{luxfi/lattice/v7} &
  Ring-LWE, NTT, sampling &
  \done & --- \\
\texttt{luxfi/node} &
  Encrypted-payload mempool admission &
  \todesign &
  \texttt{luxfi/node\#115} \\
\texttt{luxfi/consensus} &
  Block-proposer hook for deferred decryption &
  \todesign &
  \texttt{luxfi/consensus\#9} \\
\texttt{luxfi/vm} &
  Per-account FIFO mempool primitive &
  \todesign &
  \texttt{luxfi/vm\#6} \\
\texttt{luxfi/dex} &
  Decrypted-batch ingress + ordering &
  \todesign &
  \texttt{luxfi/dex\#7} \\
\texttt{luxfi/bridge} &
  FHE payload wrapping for destination-chain ingress &
  \todesign &
  \texttt{luxfi/bridge\#403} \\
\bottomrule
\end{longtable}

\noindent The honest summary: the FHE precompile, the TFHE library
(CPU + GPU NTT), the Ring-LWE primitives, and the threshold
keygen path are all production today. What is missing is the
end-to-end glue: real distributed TFHE decryption, the
encrypted-payload mempool extension, the block-proposer hook, and
the DEX / bridge integrations. Each missing piece has a tracking
issue against the relevant module.

% ================================================================
\section{Build Punch List}
% ================================================================

The build is tracked as discrete tickets across modules. Intended
ordering respects dependencies: real distributed decryption
(\texttt{luxfi/threshold\#20}) and GPU bootstrap
(\texttt{luxfi/fhe\#4}) unblock the rest; the mempool extension
(\texttt{luxfi/node\#115}) and the proposer hook
(\texttt{luxfi/consensus\#9}) unblock DEX integration
(\texttt{luxfi/dex\#7}); bridge wrapping
(\texttt{luxfi/bridge\#403}) is a follow-on once the destination
chain is encrypted-mempool capable.

\begin{enumerate}[label=\textbf{E\arabic*}]
\item \texttt{luxfi/threshold}: replace the fake-threshold TFHE
  stub with real Shamir-split-secret partial-decrypt and
  Lagrange-interpolated combine. Use the existing LSSS framework
  in \texttt{luxfi/fhe/pkg/threshold}.
\item \texttt{luxfi/fhe}: extend GPU dispatch beyond NTT to
  blind rotation, key switching, and bootstrap. Maintain
  byte-equality probe against CPU.
\item \texttt{luxfi/node}: encrypted-payload tx type; admission
  policy on signature + NIZK + bootstrap-budget bid; per-account
  FIFO partition; integration with the FHE precompile for NIZK
  verify.
\item \texttt{luxfi/consensus}: optional
  \texttt{IntentDecryptor} callback in ChainVM init; receives the
  ordered batch and returns a transformed plaintext batch within
  a configurable time budget; falls back to plaintext FIFO on
  quorum failure.
\item \texttt{luxfi/vm}: expose
  \texttt{Mempool.GetPerAccountFIFO(owner)} on the mempool
  interface; default linked.Hashmap implementation; per-account
  capacity caps.
\item \texttt{luxfi/dex}: add \texttt{BatchMetadata} (origin,
  intra-batch order) to the order ingress wire format; engine
  respects \texttt{intra\_batch\_order} for encrypted-origin
  batches.
\item \texttt{luxfi/bridge}: optional \texttt{PayloadWrapper}
  callback between signature and broadcast; resolves destination
  chain's threshold TFHE public key; encrypts under it; cancel
  path on destination decryption failure.
\end{enumerate}

% ================================================================
\section{Use Cases}
% ================================================================

The encrypted-mempool primitive is venue-neutral and substrate-level.
Downstream chains and applications consume it without modification
to the protocol described here.

\begin{itemize}[nosep]
\item \textbf{Regulated digital-securities venues} consume the
  encrypted mempool to provide best-execution guarantees to
  retail and institutional traders without an opt-in flow. The
  same defence applies uniformly to every intent.
\item \textbf{Dark pools and confidential market making} reuse
  the FHE precompile and threshold-key infrastructure to keep
  individual maker quotes private until block construction.
\item \textbf{Permissioned consortium chains} use the encrypted
  mempool to remove front-running risk on inter-bank settlement,
  even when the validator set is a small known group.
\item \textbf{Cross-chain bridges} use the bridge payload
  wrapping (\texttt{luxfi/bridge\#403}) to protect source-side
  observability of large deposits, closing the source-chain
  sandwich vector that has historically cost institutional users.
\end{itemize}

White-label applications that ship downstream of Lux import the
upstream modules and inherit the protocol without modification. The
Lux philosophy is that the primitive lives once, in the substrate,
and every consumer benefits.

% ================================================================
\section{Open Questions}
% ================================================================

\begin{enumerate}[label=\arabic*.]
\item \textbf{Throughput at peak orderbook load.} The bootstrap
  budget is the enforcement primitive; the right \emph{values}
  for per-block ops cap and per-tx bid require benchmarking at
  realistic intent volumes once the GPU-bootstrap path
  (\texttt{luxfi/fhe\#4}) lands.
\item \textbf{Threshold-key ceremony.} The threshold TFHE secret
  key is generated once across $N$ validators. Failure modes
  (partial output, share corruption, observer-in-the-middle)
  must be designed and documented before mainnet rollout.
\item \textbf{Encrypted-bridge atomicity.} An encrypted bridge
  message whose decryption fails on destination leaves a gap.
  The bridge protocol needs a retry-and-cancel path that does
  not leak content during the gap.
\item \textbf{Validator ordering attestations.} Signed
  $(\text{ciphertext-hash}, \text{arrival-time})$ tuples per
  block would let third parties audit mempool ordering without
  decryption. Design + operational cost is open.
\item \textbf{Formal mechanisation.}
  Theorems~\ref{thm:sandwich}, \ref{thm:blind-frontrun}, and
  \ref{thm:quorum-fail} are paper-form. A Lean4 or Coq
  mechanisation is planned alongside the existing
  \texttt{luxfi/proofs/} corpus.
\end{enumerate}

% ================================================================
\section{Conclusion}
% ================================================================

The Lux substrate today ships an FHE precompile, a TFHE library
with GPU NTT, a threshold-protocol framework, and a permissioned
validator option. What remains to deliver an end-to-end encrypted
mempool is precisely a glue layer: real distributed TFHE
decryption, an encrypted-payload mempool admission policy, a
proposer hook for deferred decryption, and DEX / bridge
integrations. Each piece is tracked against the relevant upstream
module. Once it lands, every consumer of the Lux stack inherits
sandwich-and-flash-loan immunity on the protected surface, with no
opt-in, no separate relay, and no per-trade configuration.

The primitive lives once. Every chain in the family, and every
white-label application downstream, gets it for free.

\end{document}
